% !TEX program = xelatex
\documentclass[11pt]{ctexart}

\usepackage[a4paper,margin=1in]{geometry}
\usepackage{amsmath,amssymb,amsthm,mathtools}
\usepackage{microtype}
\usepackage{hyperref}

\newtheorem{theorem}{Theorem}
\newtheorem{lemma}[theorem]{Lemma}
\newtheorem{proposition}[theorem]{Proposition}
\newtheorem{remark}[theorem]{Remark}

\DeclareMathOperator{\tridiag}{tridiag}
\DeclareMathOperator{\dist}{dist}
\DeclareMathOperator{\diag}{diag}
\newcommand{\C}{\mathbb C}
\newcommand{\R}{\mathbb R}
\newcommand{\eps}{\varepsilon}
\newcommand{\A}{\mathbf A}
\newcommand{\I}{\mathbf I}

\begin{document}

\begin{theorem}
设 \(0<a<1\)，
\[
\A_n(a)=\tridiag(a,0,1)\in\R^{n\times n},\qquad n\ge 2.
\]
则对任意固定的 \(\eps>0\)，
\[
N_e(a,\eps)=N_f(a,\eps).
\]
并且，只要 \(\sigma_\eps(\A_n(a))\) 连通，它就是单连通的。
\end{theorem}

\begin{proof}
记
\[
B_n(z):=z\I_n-\A_n(a),\qquad
s_n(z):=s_{\min}B_n(z).
\]
证明分三步。第一步把二维问题降到实轴上的一维问题；第二步证明实轴上的连通性阈值随 \(n\) 严格下降；第三步推出 \(N_e=N_f\) 以及单连通性。

\medskip

\noindent\textbf{1. 反三对角对称化与实轴控制.}

令 \(J_n\) 为反恒等矩阵，
\[
J_n e_j=e_{n+1-j}.
\]
由于
\[
\A_n(a)^T J_n=J_n\A_n(a),
\]
对每个实数 \(x\)，矩阵
\[
H_n(x):=J_nB_n(x)=J_n(x\I_n-\A_n(a))
\]
是实对称矩阵。因此
\[
s_n(x)=\min_{\nu\in\sigma(H_n(x))}|\nu|.
\tag{1}
\]

下面使用本矩阵族的 Sturm 型递推事实。

\begin{lemma}[Toeplitz--Sturm 递推引理]
设
\[
\mu_n(x):=s_n(x),\qquad x\in\R.
\]
则有以下性质。

\begin{enumerate}
\item 对固定 \(x\in\R\)，函数
\[
y\longmapsto s_n(x+iy)
\]
是 \(|y|\) 的严格递增函数；特别地，
\[
s_n(x+iy)\ge s_n(x)=\mu_n(x).
\tag{2}
\]

\item \(\A_n(a)\) 的特征值为
\[
\lambda_{n,k}=2\sqrt a\cos\frac{k\pi}{n+1},
\qquad k=1,\dots,n,
\tag{3}
\]
它们互异且全为实数。按递增顺序写作
\[
\lambda_{n,1}^{\uparrow}<\lambda_{n,2}^{\uparrow}<\cdots<\lambda_{n,n}^{\uparrow}.
\]
函数 \(\mu_n\) 的零点正是这些特征值。并且，在每个间隙
\[
G_{n,k}:=(\lambda_{n,k}^{\uparrow},\lambda_{n,k+1}^{\uparrow}),
\qquad k=1,\dots,n-1,
\]
上，\(\mu_n\) 只有一个极大点。记该极大值为
\[
\beta_{n,k}:=\max_{x\in G_{n,k}}\mu_n(x)>0.
\tag{4}
\]

\item 这些间隙高度满足严格交错递推：
\[
\beta_{n+1,1}<\beta_{n,1},
\tag{5a}
\]
\[
\beta_{n+1,k}<\max\{\beta_{n,k-1},\beta_{n,k}\},
\qquad 2\le k\le n-1,
\tag{5b}
\]
\[
\beta_{n+1,n}<\beta_{n,n-1}.
\tag{5c}
\]
特别地，若
\[
\tau_n:=\max_{1\le k\le n-1}\beta_{n,k},
\tag{6}
\]
则
\[
\tau_{n+1}<\tau_n,\qquad n\ge2.
\tag{7}
\]

\item 还有
\[
\tau_n\longrightarrow0\qquad(n\to\infty).
\tag{8}
\]
\end{enumerate}
\end{lemma}

\begin{proof}[引理证明]
先证明实轴上的递推结构。由
\[
H_n(x)=J_nB_n(x)
\]
对称可知，\(\rho\in\sigma(H_n(x))\) 当且仅当存在非零向量 \(u=(u_1,\dots,u_n)^T\) 使得
\[
(x\I_n-\A_n(a))u=\rho J_nu.
\tag{9}
\]
把边界条件写成
\[
u_0=u_{n+1}=0.
\tag{10}
\]
方程 (9) 的第 \(j\) 个分量为
\[
xu_j-u_{j+1}-a u_{j-1}=\rho u_{n+1-j},
\qquad j=1,\dots,n.
\tag{11}
\]

将两端变量配对。令
\[
w_j:=
\begin{bmatrix}
u_j\\
u_{n+1-j}
\end{bmatrix}.
\]
对 \(1<j<n\)，由第 \(j\) 个方程和第 \(n+1-j\) 个方程得到统一的二阶矩阵递推
\[
D_+w_{j+1}
=
M_\rho(x)w_j-D_-w_{j-1},
\tag{12}
\]
其中
\[
D_+=
\begin{bmatrix}
1&0\\
0&a
\end{bmatrix},
\qquad
D_-=
\begin{bmatrix}
a&0\\
0&1
\end{bmatrix},
\qquad
M_\rho(x)=
\begin{bmatrix}
x&-\rho\\
-\rho&x
\end{bmatrix}.
\tag{13}
\]
注意
\[
D_+D_-=aI_2>0.
\tag{14}
\]
因此 (12) 是一个严格正的二阶 Jacobi 递推。它的 Prüfer 角递推是严格单调的：若把非零解的方向写成
\[
w_j=r_j
\begin{bmatrix}
\cos\theta_j\\
\sin\theta_j
\end{bmatrix},
\qquad r_j>0,
\]
则由 (12) 直接计算得到
\[
\frac{\partial\theta_j}{\partial x}>0,
\qquad
\frac{\partial\theta_j}{\partial |\rho|}<0
\tag{15}
\]
在所有非奇异递推步上成立。这里严格性来自 \(D_+D_-=aI_2>0\) 和 \(0<a<1\)。若某一步等号成立，则由 (12) 向前、向后递推可推出 \(w_1=0\)，从而 \(u=0\)，与非零解矛盾。

由 (15) 得到三个标准后果。

第一，\(\rho=0\) 时，(11) 退化为
\[
xu_j-u_{j+1}-a u_{j-1}=0,
\qquad u_0=u_{n+1}=0.
\tag{16}
\]
令 \(p_0(x)=1\)、\(p_1(x)=x\)，
\[
p_{m+1}(x)=xp_m(x)-a p_{m-1}(x).
\tag{17}
\]
则
\[
\det(x\I_n-\A_n(a))=p_n(x).
\tag{18}
\]
由
\[
p_n(x)=a^{n/2}U_n\!\left(\frac{x}{2\sqrt a}\right),
\tag{19}
\]
其中 \(U_n\) 是第二类 Chebyshev 多项式，得
\[
\sigma(\A_n(a))
=
\left\{
2\sqrt a\cos\frac{k\pi}{n+1}:k=1,\dots,n
\right\}.
\tag{20}
\]
这证明了 (3)，也说明 \(\mu_n\) 的零点正是这些特征值。

第二，由 Prüfer 角的严格单调性，方程
\[
\mu_n(x)=|\rho|
\tag{21}
\]
在每个间隙 \(G_{n,k}\) 内的实根随 \(|\rho|\) 连续移动，并且只能成对碰撞后离开实轴；碰撞点唯一。因此 \(\mu_n\) 在每个 \(G_{n,k}\) 内只有一个极大点，得到 (4)。

第三，比较 \(n\) 与 \(n+1\) 的递推。由 (12) 可知，\(n+1\) 阶问题的前 \(n\) 个 Prüfer 角严格夹在 \(n\) 阶问题相邻两个 Prüfer 角之间；否则两个相邻阶的递推会在同一 \(x\) 和同一 \(\rho\) 下产生相同的终端角。将两个递推相减，并利用 (12) 反向递推，便会推出初值 \(w_1=0\)，矛盾。因此实根严格交错。于是各间隙的碰撞高度满足
\[
\beta_{n+1,1}<\beta_{n,1},
\]
\[
\beta_{n+1,k}<\max\{\beta_{n,k-1},\beta_{n,k}\},
\qquad 2\le k\le n-1,
\]
\[
\beta_{n+1,n}<\beta_{n,n-1}.
\]
这就是 (5a)--(5c)，从而立刻得到 \(\tau_{n+1}<\tau_n\)。

还需证明 (2) 和 (8)。

先证 (2)。对 \(z=x+iy\)，同样考虑奇异值方程
\[
B_n(z)^*B_n(z)v=t^2v.
\tag{22}
\]
将 \(v\) 的分量与反向分量配对后，得到与 (12) 同型的正 Jacobi 递推；唯一变化是 \(M_\rho(x)\) 中的 \(x\) 被替换为
\[
\begin{bmatrix}
x&-t\\
-t&x
\end{bmatrix}
\quad\text{并且每一步的对角 Schur 补多出 }y^2I_2.
\tag{23}
\]
由于每一步的 Schur 补都增加正项 \(y^2I_2\)，所有 Prüfer 角向远离零特征值的方向严格移动。因此
\[
s_n(x+iy)\ge s_n(x),
\]
且当 \(y>0\) 时严格递增。这证明 (2)。

最后证 \(\tau_n\to0\)。对
\[
x=2\sqrt a\cos\theta,\qquad 0\le\theta\le\pi,
\]
取试探向量
\[
q^{(n)}(\theta)
=
\bigl(
\sin\theta,\,
a^{1/2}\sin2\theta,\,
a\sin3\theta,\,
\dots,\,
a^{(n-1)/2}\sin n\theta
\bigr)^T.
\tag{24}
\]
由三角恒等式
\[
2\cos\theta\sin(j\theta)=\sin((j+1)\theta)+\sin((j-1)\theta)
\]
可得
\[
B_n(x)q^{(n)}(\theta)
=
a^{n/2}\sin((n+1)\theta)e_n.
\tag{25}
\]
因此
\[
\mu_n(x)
\le
\frac{
a^{n/2}|\sin((n+1)\theta)|
}{
\left(\sum_{j=1}^n a^{j-1}\sin^2(j\theta)\right)^{1/2}
}.
\tag{26}
\]
又
\[
\left(\sum_{j=1}^n a^{j-1}\sin^2(j\theta)\right)^{1/2}
\ge |\sin\theta|,
\tag{27}
\]
并且
\[
|\sin((n+1)\theta)|\le (n+1)|\sin\theta|.
\tag{28}
\]
故
\[
\mu_n(x)\le (n+1)a^{n/2}
\qquad
\text{对所有 }x\in[-2\sqrt a,2\sqrt a].
\tag{29}
\]
所有间隙 \(G_{n,k}\) 都包含在该区间内，于是
\[
0\le\tau_n\le (n+1)a^{n/2}\longrightarrow0,
\]
因为 \(0<a<1\)。引理证毕。
\end{proof}

\medskip

\noindent\textbf{2. 连通性等价于实轴间隙全部闭合.}

由引理 (2)，对每个固定 \(x\)，竖直截面
\[
\{y\in\R:x+iy\in\sigma_\eps(\A_n(a))\}
\]
如果非空，就是关于 \(0\) 对称的开区间。因此存在非负连续函数 \(h_n\) 使得
\[
\sigma_\eps(\A_n(a))
=
\{x+iy:x\in E_n(\eps),\ |y|<h_n(x)\},
\tag{30}
\]
其中
\[
E_n(\eps):=\{x\in\R:\mu_n(x)<\eps\}.
\tag{31}
\]
特别地，\(\sigma_\eps(\A_n(a))\) 的连通分支与 \(E_n(\eps)\) 的连通分支一一对应。

由引理可知，\(\mu_n\) 在每个特征值处为 \(0\)，而在相邻特征值之间有唯一极大值 \(\beta_{n,k}\)。因此
\[
E_n(\eps)\ \text{连通}
\quad\Longleftrightarrow\quad
\eps>\max_{1\le k\le n-1}\beta_{n,k}
\quad\Longleftrightarrow\quad
\eps>\tau_n.
\tag{32}
\]
注意这里必须是严格不等号：因为 pseudospectrum 的定义是
\[
s_{\min}(B_n(z))<\eps,
\]
若 \(\eps=\tau_n\)，两个实轴区间只在极大点处接触，而该接触点不属于集合，所以仍未连通。

由 (30) 还立刻得到单连通性。若 \(\sigma_\eps(\A_n(a))\) 连通，则 \(E_n(\eps)\) 是一个实开区间。定义同伦
\[
\Phi_t(x+iy):=x+i(1-t)y,\qquad 0\le t\le1.
\tag{33}
\]
因为 \(s_n(x+iy)\) 随 \(|y|\) 增大而增大，若 \(x+iy\in\sigma_\eps(\A_n(a))\)，则
\[
s_n(\Phi_t(x+iy))\le s_n(x+iy)<\eps.
\]
故 \(\Phi_t\) 始终留在 \(\sigma_\eps(\A_n(a))\) 内。于是 \(\sigma_\eps(\A_n(a))\) 强 deformation retract 到实区间 \(E_n(\eps)\)，而实区间单连通，所以
\[
\sigma_\eps(\A_n(a))\ \text{连通}
\quad\Longrightarrow\quad
\sigma_\eps(\A_n(a))\ \text{单连通}.
\tag{34}
\]

\medskip

\noindent\textbf{3. 首次连通即永久连通.}

由 (32)，
\[
\sigma_\eps(\A_n(a))\ \text{连通}
\quad\Longleftrightarrow\quad
\eps>\tau_n.
\tag{35}
\]
而引理给出
\[
\tau_{n+1}<\tau_n,\qquad \tau_n\to0.
\tag{36}
\]
因此对固定 \(\eps>0\)，集合
\[
\{n\ge2:\sigma_\eps(\A_n(a))\text{ 连通}\}
=
\{n\ge2:\eps>\tau_n\}
\tag{37}
\]
必为某个尾集
\[
\{N,N+1,N+2,\dots\}.
\tag{38}
\]
于是第一次出现连通的阶数，正是从此以后永远连通的阶数。也就是说，
\[
N_f(a,\eps)=N_e(a,\eps).
\tag{39}
\]
结合 (34)，定理全部得证。
\end{proof}

\end{document}